% riess_schmidt_perlmutter_ML.tex
% "The Accelerating Universe as Statistical Inference:
%  Type Ia Supernovae, Maximum Likelihood, and the
%  Cosmological Constant, 1988--1999"
% Thomas J. Sargent, 2026

\documentclass[12pt]{article}

\usepackage{amsmath,amssymb}
\usepackage{booktabs}
\usepackage{array}
\usepackage{multirow}
\usepackage[margin=1.2in]{geometry}
\usepackage{setspace}
\usepackage{microtype}
\usepackage{hyperref}
\usepackage{natbib}
\usepackage{parskip}

\onehalfspacing

\hypersetup{
  colorlinks  = true,
  linkcolor   = blue,
  citecolor   = blue,
  urlcolor    = blue
}

\title{The Accelerating Universe as Statistical Inference:\\
       Type Ia Supernovae, Maximum Likelihood,\\
       and the Cosmological Constant, 1988--1999}
\author{Thomas J.\ Sargent\thanks{%
  In 2011 my good friend Christopher A.\ Sims and I shared the Nobel Prize in
  economics.  That good luck let us meet and become friends with winners of
  prizes in physics and other fields.  Our conversations with the physics
  winners were especially interesting.  Adam Riess had majored partly in
  economics as an undergraduate and told us about his very interesting
  undergraduate thesis.  I have kept in contact with Adam Riess and Brian
  Schmidt, who have taught me the nuts and bolts of the computer programs that
  they used to compute and climb the likelihood function that led to their
  remarkable discovery.
  This paper is one of five that read episodes in the
  history of physics and astronomy as exercises in statistical inference and
  machine learning.  \citet{sargent2025sources} sets out the theme: what is
  now called artificial intelligence names practices that scientists have used
  for centuries.  The other four treat Kepler and Newton, Einstein's photoelectric equation,
  Pauli's exclusion principle, and the exoplanet discovery of Mayor and
  Queloz.}}
\date{April 2026}

\begin{document}
\maketitle

\begin{abstract}
In their Nobel Prize--winning papers, Adam Riess, Brian Schmidt, and their
High-Z Supernova Search Team \citep{Riess1998}, and independently Saul
Perlmutter and the Supernova Cosmology Project \citep{Perlmutter1999},
reported that the expansion of the universe is \emph{accelerating}.
Their argument had two components.  A physical \emph{model} supplied the
Friedmann equations with a cosmological constant $\Lambda$.  A
\emph{measurement equation} related the apparent magnitude of a Type Ia
supernova to its redshift, with the shape of that relation controlled by two
density parameters~$\Omega_M$ and~$\Omega_\Lambda$.  This note examines their
reasoning from a modern statistical perspective.  I describe the data both
groups assembled, the Cal\'an/Tololo low-redshift training sample and the
high-redshift discoveries at $z = 0.18$--$0.83$.  I then lay out the maximum
likelihood problem they solved: choose $(\Omega_M, \Omega_\Lambda)$ to
minimise the weighted sum of squared residuals between observed distance
moduli and the Friedmann prediction.  Under Gaussian errors MLE coincides with
weighted least squares.  The curvature of the log-likelihood surface---the Fisher
information---determines the confidence ellipses in the
$(\Omega_M, \Omega_\Lambda)$ plane.  A nuisance parameter, the combination
$\mathcal{M}$ of the absolute magnitude~$M_B$ and the Hubble
constant~$H_0$, enters the measurement equation as the contact
electromotive force enters Millikan's photoelectric regression.  It shifts
the Hubble diagram vertically and leaves the shape of $\mu(z)$ alone, so the
curvature of the diagram identifies the shape parameters independently
of~$\mathcal{M}$.  The analogy stops there.  Millikan's slope was identified;
these two parameters are identified only in one combination.  The estimates
under flatness, $\Omega_M \approx 0.28$ and
$\Omega_\Lambda \approx 0.72$, imply a deceleration parameter
$q_0 = \Omega_M/2 - \Omega_\Lambda \approx -0.58$.

I compute the Fisher information for the 1998 designs.  It reproduces the
published error bar on the combination $0.8\Omega_M - 0.6\Omega_\Lambda$,
$\pm0.1$, from the redshift distributions and quoted uncertainties alone.  It
also shows that the individual parameters are not measured: the marginal
standard errors are $0.44$ and $0.65$.  The published $\pm0.09$ is the
standard error under the flatness constraint.  Ruling out
$\Omega_\Lambda = 0$ therefore requires a prior, either $\Omega_M \geq 0$ or
flatness from the cosmic microwave background.

I read the episode in the language of machine learning.  The Friedmann
equations supplied the hypothesis class.  The Phillips width--luminosity
relation supplied the feature that made the signal visible, as Lenard's
finding that frequency rather than intensity governs photoemission supplied
Einstein's.  A deep network given the same data recovers neither the density
parameters nor their meaning.  The supernovae measure one function,
$d_L(z)$, and every theory reproducing it remains standing: theory reduces the
sample size that data must supply, and theory also names the measurement that
discriminates.
\end{abstract}
%
%\tableofcontents
%
%======================================================================
\section{Introduction}

In 1998 two independent teams reported that distant Type Ia supernovae are
systematically \emph{fainter} than they would be in a universe whose
expansion is decelerating.  \citet{Riess1998} appeared in September 1998 and
\citet{Perlmutter1999} in June 1999; both groups had announced preliminary
results at meetings earlier in 1998.  The simplest explanation is that the
expansion is speeding up, driven by a positive cosmological
constant~$\Lambda$ or some analogous dark energy.  Saul Perlmutter received
one half of the 2011 Nobel Prize in Physics for the discovery, and Adam Riess
and Brian Schmidt shared the other half.

A twenty-first-century statistician confronted with tables of
(redshift, apparent magnitude) pairs would recognise the problem as a
weighted nonlinear regression with two scientifically meaningful
parameters and a nuisance intercept.  The Friedmann equations with a
cosmological constant were well established before the data were collected.
The discovery consisted in determining the sign and magnitude of
$\Omega_\Lambda$.

Six features of the episode interest a statistician.
\begin{enumerate}
  \item \textbf{Model specification: the hypothesis class from the
    Friedmann equations.}  General relativity predicts the
    luminosity-distance--redshift relation
    $d_L(z;\Omega_M,\Omega_\Lambda)$ exactly.  The functional form is
    fixed, and the estimation problem reduces to two dimensionless density
    parameters.
  \item \textbf{Standardising a noisy standard candle.}  The peak $B$-band
    luminosities of raw Type~Ia supernovae scatter by
    $\approx 0.4$~magnitudes.  \citet{Phillips1993} found that the rate of
    post-peak decline~$\Delta m_{15}$ predicts peak luminosity, which
    reduces the dispersion to $\approx 0.15$~mag.
  \item \textbf{The nuisance parameter~$\mathcal{M}$.}  The combination of
    absolute magnitude~$M_B$ and Hubble constant~$H_0$ shifts the entire
    Hubble diagram vertically.  The observations do not identify
    $\mathcal{M}$ on its own, and both teams profiled it out of the
    likelihood.  The relative distance moduli of high-$z$ and low-$z$
    supernovae identify the shape of $\mu(z)$.
  \item \textbf{Partial identification of the density parameters.}  The
    shape of $\mu(z)$ pins down one combination of $\Omega_M$ and
    $\Omega_\Lambda$ and leaves the orthogonal combination nearly free.  The
    published error bars assume flatness.  Acceleration follows from the
    supernova data together with a prior, not from the supernova data
    alone.
  \item \textbf{Independent replication with different methods.}
    The High-Z Supernova Search Team used the MLCS
    (Multi-Light-Curve Shape) standardisation method of
    \citet{RiessPress1996}; the Supernova Cosmology Project used the
    stretch-factor method of \citet{Perlmutter1997}.  The two groups
    analysed different supernovae and reached the same conclusion to
    within their measurement uncertainties.
  \item \textbf{Reluctance to accept the model.}  Many physicists doubted
    $\Omega_\Lambda > 0$ because the prevailing theory held that the vacuum
    energy density should be negligible \citep{Weinberg1989}.
\end{enumerate}

Both teams had roughly 40--60 high-quality supernovae, a meagre dataset by
modern standards.  The Friedmann equations specified the regression function
up to three free parameters and so collapsed a high-dimensional
function-approximation task into a low-dimensional one.  Even that proved too
large for the data.  Fifty supernovae determine one combination of the two
density parameters and leave the other open, and the published estimates close
it with a prior.

Sections~\ref{sec:background}--\ref{sec:data} set the observational and
theoretical stage.
Section~\ref{sec:wls} presents the WLS/MLE procedure.
Section~\ref{sec:fisher} develops the Fisher information and confidence
ellipses.
Section~\ref{sec:nuisance} discusses the $\mathcal{M}$ identification
problem.
Section~\ref{sec:results} presents the main results and the inference
of acceleration.
Section~\ref{sec:twogroups} compares the two teams.
Section~\ref{sec:reluctance} describes the scepticism that greeted the
result and the robustness tests that answered it.
Section~\ref{sec:prediction} describes the prediction that supernovae at
$z > 1$ later confirmed.
Section~\ref{sec:limits} describes what the supernovae could not settle.
Sections~\ref{sec:ml}--\ref{sec:dnn} reinterpret the episode in the
language of machine learning.
Section~\ref{sec:conclusions} concludes.

A companion paper treats Millikan's 1916 photoelectric experiment in the same
way, and the comparison runs through what follows.  Einstein's equation is a
line whose slope is a universal constant, fixed in advance by Planck's
blackbody fit and predicted with no free parameter.  A contact electromotive
force shifted every reading by the same amount and left the slope untouched,
as $\mathcal{M}$ shifts every point of the Hubble diagram.  Millikan measured
the slope with five observations to better than one percent.  The disanalogy
matters more than the analogy: Millikan had one structural parameter, so
removing the nuisance left it identified, while the Friedmann model has two
and removing the nuisance leaves a second degeneracy behind.
Section~\ref{sec:fisher} turns on that difference.

%======================================================================
\section{Observational Background, 1929--1994}
\label{sec:background}

\subsection{Hubble's Law and the Standard Model of Cosmology}

In 1929 Edwin Hubble \citeyearpar{Hubble1929} reported that the
recession velocities~$v$ of nearby ``extra-galactic nebulae'' satisfy
\begin{equation}\label{eq:hubble}
  v = H_0\,d,
\end{equation}
where $d$ is the distance and $H_0$ the Hubble constant.  Hubble
estimated $H_0 \approx 500$\,km\,s$^{-1}$\,Mpc$^{-1}$; subsequent
Cepheid-variable distance rescaling reduced the accepted value to
$H_0 \approx 70$\,km\,s$^{-1}$\,Mpc$^{-1}$.  Alexander Friedmann had
already shown in 1922 and 1924 \citep{Friedmann1922,Friedmann1924} that
expanding cosmological solutions follow from Einstein's field equations;
Hubble's measurement confirmed that we live in one of them.

In the Friedmann model, the dynamical variable is the scale factor
$a(t)$, normalised so that $a(t_0) = 1$ today.  The recession redshift
$z$ of a source satisfies $1+z = 1/a(t_{\rm emit})$.  Three dimensionless
density parameters characterise the matter-energy content:
\begin{align}
  \Omega_M &= \frac{8\pi G\rho_{M,0}}{3H_0^2},
  &\text{(non-relativistic matter)}\label{eq:omegaM}\\[4pt]
  \Omega_\Lambda &= \frac{\Lambda}{3H_0^2},
  &\text{(cosmological constant / dark energy)}\label{eq:omegaL}\\[4pt]
  \Omega_k &= \frac{-k}{a_0^2 H_0^2},
  &\text{(spatial curvature)}\label{eq:omegak}
\end{align}
with the constraint $\Omega_M + \Omega_\Lambda + \Omega_k = 1$.  A flat
universe has $\Omega_k = 0$, i.e., $\Omega_M + \Omega_\Lambda = 1$.

Einstein had introduced $\Lambda > 0$ in 1917 \citeyearpar{Einstein1917}
to obtain a static universe, then abandoned it after Hubble established
expansion.  The two teams' discovery that $\Omega_\Lambda \approx 0.72$
rehabilitated the cosmological constant as an observational reality.

\subsection{Type Ia Supernovae as Standard Candles}
\label{sec:standard_candle}

A \emph{standard candle} is an astronomical object of known intrinsic
luminosity.  If absolute luminosity $L$ is known, the observed flux
$F = L/(4\pi d_L^2)$ gives the luminosity distance $d_L$ directly.
Comparing $d_L$ with the redshift $z$ maps the distance-redshift relation
and thereby constrains the Friedmann parameters.

Type Ia supernovae are thermonuclear explosions of carbon-oxygen white
dwarfs in binary systems \citep{Filippenko1997}.  The account current in the
1990s had the white dwarf accreting until it approached the Chandrasekhar
limit $M_{\rm Ch} \approx 1.4\,M_\odot$, so that a near-universal exploding
mass produced a near-universal luminosity.  That account no longer commands
consensus.  Sub-Chandrasekhar double detonations and white-dwarf mergers are
now taken seriously, and the progenitor population may well be mixed.

The uniformity of Type~Ia peak luminosity is therefore an empirical fact
rather than a derived one, and it was empirical in 1998 too.  Before
standardisation the scatter in peak $B$-band absolute magnitude is
$\sigma_{M_B} \approx 0.4$~mag.  Since five magnitudes is a factor of one
hundred in flux, that is a $20\%$ distance uncertainty, too large for
cosmology.

\subsection{The Width-Luminosity Relation: The Key Standardisation Insight}

\citet{Phillips1993} established a tight empirical relation between the
peak luminosity of a Type Ia supernova and the speed at which its light
curve declines after maximum light.  The parameter $\Delta m_{15}(B)$
measures the drop in $B$-band magnitude during the 15 days after peak:
brighter (more energetic) supernovae have slower, broader light curves
($\Delta m_{15}$ smaller), while fainter events decline more rapidly.
Quantitatively, the relation is approximately
\begin{equation}\label{eq:phillips}
  M_B^{\rm peak} \approx M_B^{\rm ref} + \alpha\bigl(\Delta m_{15} - 1.1\bigr),
  \quad \alpha \approx 0.75~\text{mag},
\end{equation}
where $M_B^{\rm ref}$ is the absolute magnitude of a ``normal'' event with
$\Delta m_{15} = 1.1$.  The coefficient $\alpha \approx 0.75$ is the
\citet{Hamuy1996} calibration.  \citet{Phillips1993} himself fitted a
steeper slope, near $2.7$~mag per unit $\Delta m_{15}(B)$, to nine
supernovae; later samples flattened it.  Both figures come from secondary
sources and should be checked against the originals.

After the correction the intrinsic dispersion falls to
$\sigma_{\rm int} \approx 0.15$~mag.  Published values run from $0.13$ to
$0.17$~mag depending on sample and method, and the two teams quote $0.13$ and
$0.14$.  I use $0.15$~mag throughout for the intrinsic dispersion and
$0.22$~mag for the total per-supernova error at $z \approx 0.5$, which adds
photometric and $K$-correction terms.

Without the correction the scatter masks the cosmological signal.  With it
the $\mu(z)$ residuals separate a decelerating from an accelerating universe.

The Phillips relation is often called the supernova counterpart of Lenard's
discovery that the photoelectric stopping potential depends on frequency and
not intensity.  The two findings did different work.  Frequency and intensity
were both physical primitives, and Lenard's experiment chose between
candidates already on the table.  Nothing put $\Delta m_{15}$ on the table.
Phillips found it by fitting curves to light curves, no theory predicted it
beforehand, and none derives it now.  Section~\ref{sec:phillips_induction}
takes up what that difference means for the reading of the episode offered
here.

\subsection{The Two Teams}

The \textbf{Supernova Cosmology Project} (SCP), led by Saul Perlmutter
at Lawrence Berkeley National Laboratory, began systematic high-$z$
supernova searches in the late 1980s and developed a systematic discovery
pipeline using ground-based telescopes (Isaac Newton Telescope, Cerro
Tololo Inter-American Observatory, Keck).
\citet{GoobarPerlmutter1995} describe their ``batch mode'' scheduling.  The
team discovered candidates in a first observing run and followed them up in a
prescheduled second run three weeks later.  The schedule guaranteed that the
events were still near peak brightness when the spectroscopy was taken.  By
1997 they had reported results for seven high-$z$ supernovae
\citep{Perlmutter1997}.  Their definitive 42-supernova paper appeared in 1999
\citep{Perlmutter1999}.

The \textbf{High-Z Supernova Search Team} (HZT), assembled around
Brian Schmidt at the Mount Stromlo Observatory and Robert Kirshner at
Harvard, made complementary use of the Hubble Space Telescope and
larger ground-based facilities to obtain high-quality light curves and
spectra.  Adam Riess developed the MLCS analysis pipeline for the team's
data \citep{RiessPress1996}.  A first report on four HST supernovae
\citep{Garnavich1998} was followed by the definitive 16-plus-34 SN paper
of \citet{Riess1998}.

%======================================================================
\section{The Cosmological Model and the Measurement Equation}
\label{sec:model}

\subsection{The Friedmann Equation}

The first Friedmann equation with matter density~$\rho_M$, pressure-free
dust, and cosmological constant~$\Lambda$ is
\begin{equation}\label{eq:friedmann}
  H^2(z) \equiv \left(\frac{\dot a}{a}\right)^2
  = H_0^2\Bigl[\Omega_M(1+z)^3 + \Omega_k(1+z)^2 + \Omega_\Lambda\Bigr]
  \equiv H_0^2\,E^2(z),
\end{equation}
where $z$ is evaluated on the past light cone and the present-day
normalisation $E(0) = 1$ implies $\Omega_M + \Omega_k + \Omega_\Lambda = 1$.
The dimensionless Hubble function $E(z)$ gives the expansion rate at
redshift $z$ in units of $H_0$.  For $\Omega_\Lambda = 0$ and matter-only
content, $E(z) = (1+z)^{3/2}$ (the Einstein-de Sitter solution); for
the best-fit flat $\Lambda$CDM model, $E(z)$ passes through intermediate
values between these extremes.

\subsection{Luminosity Distance}

The luminosity distance to a source at redshift $z$ is
\begin{equation}\label{eq:dl}
  d_L(z;\Omega_M,\Omega_\Lambda) =
  \frac{c(1+z)}{H_0\sqrt{|\Omega_k|}}
  \,f\!\left(\sqrt{|\Omega_k|}\int_0^z\frac{dz'}{E(z')}\right),
\end{equation}
where
\begin{equation}
  f(x) =
  \begin{cases}
    \sin x   & \Omega_k < 0 \text{ (closed)}, \\
    x         & \Omega_k = 0 \text{ (flat)}, \\
    \sinh x  & \Omega_k > 0 \text{ (open)}.
  \end{cases}
\end{equation}
For the flat case of primary interest ($\Omega_k = 0$, $\Omega_M +
\Omega_\Lambda = 1$), equation~\eqref{eq:dl} reduces to
\begin{equation}\label{eq:dl_flat}
  \boxed{d_L(z;\Omega_M) =
  \frac{c(1+z)}{H_0}\int_0^z\frac{dz'}
  {\sqrt{\Omega_M(1+z')^3 + (1-\Omega_M)}}.}
\end{equation}
Here $\Omega_\Lambda = 1 - \Omega_M$ is not a free parameter once flatness
is imposed.  For the general (non-flat) case, $\Omega_M$ and
$\Omega_\Lambda$ are estimated jointly without the flatness constraint.

\subsection{The Distance Modulus and Measurement Equation}

The \emph{distance modulus} is defined as
\begin{equation}\label{eq:mu_def}
  \mu \equiv m_B - M_B = 5\log_{10}\!\left(\frac{d_L}{10~\text{pc}}\right),
\end{equation}
where $m_B$ is the observed peak apparent magnitude (standardised for
light-curve shape) and $M_B$ is the intrinsic absolute magnitude.  Since
$1~\text{Mpc} = 10^6~\text{pc}$, this becomes
\begin{equation}\label{eq:mu_mpc}
  \mu = 5\log_{10}\!\left(\frac{d_L}{\text{Mpc}}\right) + 25.
\end{equation}
The theoretical distance modulus predicted by the cosmological model is
therefore
\begin{equation}\label{eq:mu_theory}
  \mu^{\rm th}(z;\Omega_M,\Omega_\Lambda,H_0) =
  5\log_{10}\!\left(\frac{d_L(z;\Omega_M,\Omega_\Lambda)}{\text{Mpc}}\right)
  + 25.
\end{equation}
The full measurement equation---the fundamental regression of the entire
analysis---is
\begin{equation}\label{eq:meas}
  \boxed{m_{B,i} = M_B + \mu^{\rm th}(z_i;\Omega_M,\Omega_\Lambda,H_0)
    + \varepsilon_i,
    \quad \varepsilon_i \sim \mathcal{N}(0,\sigma_i^2),}
\end{equation}
The sum of squared errors defines the likelihood.  Four quantities appear on
the right, $M_B$, $H_0$, $\Omega_M$ and $\Omega_\Lambda$, but only three
combinations of them are estimable, as the next subsection shows.

\subsection{The Nuisance Parameter $\mathcal{M}$}
\label{sec:mathcal_M}

Expanding equation~\eqref{eq:mu_theory} reveals that $H_0$ and $M_B$
enter only through a single combination.  Write
\begin{equation}\label{eq:dl_split}
  d_L(z;\Omega_M,\Omega_\Lambda,H_0)
  = \frac{c}{H_0}\,\tilde d_L(z;\Omega_M,\Omega_\Lambda),
\end{equation}
where $\tilde d_L$ is the \emph{dimensionless} luminosity distance that
depends only on the density parameters.  Then
\begin{align}
  \mu^{\rm th} &= 5\log_{10}\!\left(\frac{c}{H_0}\,
    \tilde d_L(z;\Omega_M,\Omega_\Lambda)\right) + 25 \notag \\
  &= 5\log_{10}\tilde d_L(z;\Omega_M,\Omega_\Lambda)
    + 5\log_{10}\!\left(\frac{c}{H_0\,\text{Mpc}}\right) + 25,
\end{align}
so the measurement equation~\eqref{eq:meas} becomes
\begin{equation}\label{eq:meas2}
  m_{B,i} = \underbrace{M_B + 5\log_{10}\!\left(\frac{c}{H_0\,\text{Mpc}}\right) + 25}_{
    \displaystyle\mathcal{M}}
    + 5\log_{10}\tilde d_L(z_i;\Omega_M,\Omega_\Lambda) + \varepsilon_i.
\end{equation}
The scalar $\mathcal{M}$ shifts the Hubble diagram vertically; the term
$5\log_{10}\tilde d_L(z;\Omega_M,\Omega_\Lambda)$ determines its shape and
does not depend on $\mathcal{M}$.
\begin{itemize}
  \item $\mathcal{M}$ is not identified by the shape of the Hubble
    diagram.  Any vertical shift of the entire $m_B$ sequence can be
    absorbed into $\mathcal{M}$.
  \item The \emph{curvature} of $\mu(z)$---the way the slope of the
    $m_B$--$z$ diagram changes between low-$z$ and high-$z$ supernovae---is
    informative about $\Omega_M$ and $\Omega_\Lambda$, and carries no
    information about $\mathcal{M}$.
\end{itemize}
Curvature separates the shape parameters from $\mathcal{M}$.  It does not
separate them from each other.  Over the observed redshift range a rise in
$\Omega_M$ and a rise in $\Omega_\Lambda$ bend $\mu(z)$ in nearly opposite
ways, so the data fix one combination of the two and leave the orthogonal
combination nearly free.  Section~\ref{sec:fisher} quantifies this, and
Section~\ref{sec:nuisance} compares the identification structure with
Millikan's contact-EMF problem.

The estimating model has three effective parameters:
\begin{equation}\label{eq:param_set}
  \boldsymbol\theta = (\mathcal{M},\,\Omega_M,\,\Omega_\Lambda),
\end{equation}
plus an intrinsic scatter~$\sigma_{\rm int}$.  For a flat universe the
constraint $\Omega_M + \Omega_\Lambda = 1$ reduces the shape parameters
to the single free parameter~$\Omega_M$.

%======================================================================
\section{The Historical Data}
\label{sec:data}

\subsection{The Low-Redshift Calibration Sample}

Measuring $(\Omega_M,\Omega_\Lambda)$ requires supernovae spanning a
wide enough redshift range that the \emph{shape} of the $\mu(z)$
curve---which changes between different cosmological models---is
statistically distinguishable.  Both teams anchored their analyses with
a common low-redshift ($z \lesssim 0.1$) training set: the
Cal\'an/Tololo Supernova Survey of \citet{Hamuy1996}, which provided 29
carefully observed Type Ia supernovae in $BVI$ bands.  The SCP used 18 of
them; the HZT drew its 34 low-$z$ events more widely from the literature.
Table~\ref{tab:calantololo} lists a representative subset.

\begin{table}[ht]
\centering
\caption{Representative Type~Ia supernovae from the Cal\'an/Tololo
low-redshift sample \citep{Hamuy1996}.  The effective $B$-band peak
magnitude $m_B^{\rm eff}$ incorporates the light-curve shape correction
of equation~\eqref{eq:phillips}.  These events serve as the
``training data'' in section~\ref{sec:ml}: they calibrate the absolute
magnitude~$M_B$ for use in the high-$z$ sample.}
\label{tab:calantololo}
\medskip
\begin{tabular}{lrrc}
\toprule
Supernova & $z$ & $m_B^{\rm eff}$ (mag) & $\sigma$ (mag) \\
\midrule
SN 1990O  & 0.030 & 16.33 & 0.07 \\
SN 1992P  & 0.026 & 16.08 & 0.07 \\
SN 1992ae & 0.075 & 18.44 & 0.09 \\
SN 1992al & 0.014 & 14.60 & 0.06 \\
SN 1992bc & 0.020 & 15.18 & 0.06 \\
SN 1992bh & 0.045 & 17.42 & 0.08 \\
SN 1992bl & 0.043 & 17.29 & 0.08 \\
SN 1992bo & 0.018 & 15.95 & 0.06 \\
SN 1993ag & 0.050 & 17.68 & 0.09 \\
SN 1993O  & 0.052 & 17.72 & 0.08 \\
\bottomrule
\end{tabular}
\end{table}

At low redshifts $d_L \approx cz/H_0$ (Hubble's law), so
$\mu \approx 5\log_{10}(cz/H_0\cdot\text{Mpc}) + 25$.  The shape parameters
enter at order $z$, through the combination $q_0 = \Omega_M/2 -
\Omega_\Lambda$, and Appendix~\ref{app:lowz} works out the coefficient.  They
are nonetheless unmeasurable from $z \lesssim 0.1$ data, because the leverage
$\sum_i(z_i-\bar z)^2$ over so narrow a range is too small: the low-$z$ sample
alone determines $q_0$ only to $\pm 1.1$.  It constrains $\mathcal{M}$ and
leaves the cosmology open.

\subsection{The High-Redshift Samples}

Both teams spent several years assembling high-$z$ supernova samples
using ground-based telescopes supplemented by Hubble Space Telescope
photometry and Keck spectroscopy.

\paragraph{Perlmutter et al.\ (1999).}  The SCP's ``Fit~C'' sample
comprises 18 high-quality high-$z$ events; their full ``Fit~C+D''
sample extends to 42 supernovae at $z = 0.18$--$0.83$.  Table
\ref{tab:hiZ_scp} lists a representative subset, showing the effective
apparent magnitudes after $K$-correction (accounting for the redshift
of the SN Ia spectral energy distribution relative to the $B$-band
filter) and stretch-factor standardisation.

\begin{table}[ht]
\centering
\caption{Representative high-redshift Type~Ia supernovae from the
Supernova Cosmology Project \citep{Perlmutter1999}.  Effective
$B$-band magnitudes~$m_B^{\rm eff}$ after stretch-factor and
$K$-correction.  The photometric uncertainty~$\sigma$ includes
contributions from photon noise, host-galaxy subtraction, $K$-correction
uncertainty, and intrinsic scatter~$\sigma_{\rm int} \approx 0.15$~mag.
Values are approximately reproduced from Table~1 of
\citet{Perlmutter1999}.}
\label{tab:hiZ_scp}
\medskip
\begin{tabular}{lrrc}
\toprule
Supernova & $z$ & $m_B^{\rm eff}$ (mag) & $\sigma$ (mag) \\
\midrule
SN 1994am & 0.372 & 22.46 & 0.21 \\
SN 1994H  & 0.374 & 22.17 & 0.15 \\
SN 1994an & 0.378 & 22.92 & 0.18 \\
SN 1995az & 0.450 & 22.79 & 0.20 \\
SN 1994al & 0.420 & 22.45 & 0.22 \\
SN 1994G  & 0.425 & 22.55 & 0.17 \\
SN 1994F  & 0.354 & 22.24 & 0.19 \\
SN 1995ba & 0.388 & 22.65 & 0.17 \\
SN 1992bi & 0.458 & 22.92 & 0.16 \\
SN 1996cg & 0.490 & 23.38 & 0.18 \\
SN 1997ap & 0.830 & 24.23 & 0.23 \\
\bottomrule
\end{tabular}
\end{table}

\paragraph{Riess et al.\ (1998).}  The HZT's primary dataset comprised
16 high-$z$ supernovae ($z = 0.16$--$0.97$), supplemented by 34
low-$z$ events compiled from the literature.  The MLCS method was
applied to multi-band ($BVRI$) light curves, providing simultaneous
estimates of the shape parameter~$\Delta$, the colour excess
$E(B-V)$ (dust reddening), and the peak magnitude.  A subset is
shown in Table~\ref{tab:hiZ_hzt}.

\begin{table}[ht]
\centering
\caption{Selected high-redshift supernovae from the High-Z Supernova
Search Team \citep{Riess1998}.  Corrected distance moduli $\mu_{\rm MLCS}$
are computed from the full MLCS fit to the multi-band light curves
and incorporate corrections for light-curve shape ($\Delta$) and
dust reddening ($A_B$).  The MLCS zero point is arbitrary and is absorbed
into $\mathcal{M}$, so these moduli carry a common offset relative to any
particular cosmology; only their differences matter.  Values approximately
reproduced from \citet{Riess1998}, Table~5.}
\label{tab:hiZ_hzt}
\medskip
\begin{tabular}{lrrrc}
\toprule
Supernova & $z$ & $\mu_{\rm MLCS}$ & $\sigma_\mu$ & Notes \\
\midrule
SN 1997ce & 0.440 & 42.53 & 0.33 & ``golden'' \\
SN 1997cj & 0.500 & 42.91 & 0.24 & ``golden'' \\
SN 1997ck & 0.970 & 44.30 & 0.51 & no spectrum \\
SN 1995ao & 0.300 & 41.08 & 0.36 & \\
SN 1995ap & 0.230 & 40.36 & 0.26 & \\
SN 1996cf & 0.570 & 43.16 & 0.42 & \\
SN 1996bj & 0.574 & 43.67 & 0.26 & \\
\bottomrule
\end{tabular}
\end{table}

\subsection{Provenance of the Tabulated Data}
\label{sec:provenance}

Tables~\ref{tab:calantololo}--\ref{tab:hiZ_hzt} are representative subsets
transcribed from the published tables.  I have not checked every entry against
the originals, and no result in this paper is computed from them.  The
parameter estimates in Section~\ref{sec:results} are the two teams' own
published values.  The forecasts in Section~\ref{sec:fisher} use the observed
redshift distributions and quoted uncertainties, not the individual magnitudes.

The tables are still informative about scale.  Referred to a flat $\Lambda$CDM
model with $H_0 = 70$\,km\,s$^{-1}$\,Mpc$^{-1}$, the Cal\'an/Tololo entries
imply $M_B$ between $-18.5$ and $-19.5$, with a mean near $-19.2$ and a
dispersion of about $0.25$~mag.  The SCP high-$z$ entries imply a mean $M_B$
near $-19.1$.  The difference between the two means is the acceleration signal.
The difference is about $0.1$~mag in this subset, against the $0.25$~mag that
the full samples deliver.  A dozen events with $0.2$~mag errors show no more.

\subsection{The Standardisation Pipeline}

Both teams corrected the raw observed magnitudes for four effects
before entering them into the cosmological likelihood:

\begin{enumerate}
  \item \textbf{$K$-correction.}  At $z \sim 0.5$, the observer-frame
    $I$-band covers rest-frame $B$-band, while the observer-frame
    $B$-band covers rest-frame UV.  A wavelength-dependent correction
    transforms the measured filter magnitude to the rest-frame $B$-band
    magnitude that anchors the low-$z$ calibration.  The SCP used
    templates of \citet{Kim1996}; the HZT used spectrophotometric
    templates from \citet{Nugent2002}.

  \item \textbf{Light-curve shape correction.}  The SCP employed a
    \emph{stretch factor}~$s$ (time-axis scaling): each light curve is
    remapped to the template by stretching the time axis by $s$, then
    the peak magnitude is corrected by $\alpha_{\rm SCP}(s-1)$ with
    $\alpha_{\rm SCP}\approx 1.35$~mag.  The HZT's MLCS method fit
    multi-band templates parameterised by a single luminosity
    offset~$\Delta$; the correction is~$\eta\,\Delta$ with
    $\eta \approx 0.91$~mag per unit $\Delta$.

  \item \textbf{Colour/dust correction.}  Both teams corrected for
    dust reddening in the host galaxy using the measured $B-V$ colour
    excess at peak light: $A_B = R_B\,E(B{-}V)$ with
    $R_B \approx 4.1$.

  \item \textbf{Peculiar-velocity uncertainty.}  Each SN at $z < 0.1$
    contributes a velocity uncertainty $\sigma_v \approx 200$--$300$\,km/s
    from large-scale structure motions, adding an effective distance-
    modulus error $\sigma_{v,i} = 5\sigma_v/(cz_i \ln 10)$ that diverges
    at low $z$.  This is why the minimum useful redshift is $z \approx 0.01$.
\end{enumerate}
After these corrections, the effective $m_B^{\rm eff}$ enters the
cosmological likelihood as if it were the true distance modulus plus a
Gaussian residual.

The total photometric uncertainty per supernova is typically
\[
  \sigma_{m,i}^2 = \sigma_{\rm photo,i}^2
    + \sigma_{K,i}^2 + \sigma_{v,i}^2 + \sigma_{\rm int}^2,
\]
where $\sigma_{\rm int}$ is the residual dispersion after standardisation,
taken as $0.15$~mag throughout (Section~\ref{sec:standard_candle}).

%======================================================================
\section{Weighted Least Squares and Maximum Likelihood Estimation}
\label{sec:wls}

\subsection{The Likelihood Function}

Given the measurement model~\eqref{eq:meas2}, in which each observation
$m_{B,i}$ is drawn from a normal distribution centred on $\mathcal{M} +
5\log_{10}\tilde d_L(z_i;\Omega_M,\Omega_\Lambda)$, the likelihood for
a sample of $n = N_{\rm low} + N_{\rm high}$ supernovae is
\begin{equation}\label{eq:lik}
  L(\boldsymbol\theta) =
  \prod_{i=1}^n \frac{1}{\sqrt{2\pi}\,\sigma_i}
  \exp\!\left(-\frac{\bigl[m_{B,i} - \mathcal{M}
    - 5\log_{10}\tilde d_L(z_i;\Omega_M,\Omega_\Lambda)\bigr]^2}
    {2\sigma_i^2}\right),
\end{equation}
where the total uncertainty $\sigma_i$ (including intrinsic scatter)
was discussed in Section~\ref{sec:data}.  The log-likelihood is
\begin{equation}\label{eq:loglik}
  \ell(\boldsymbol\theta) =
  -\frac{n}{2}\ln(2\pi)
  - \sum_{i=1}^n \ln\sigma_i
  - \underbrace{\frac{1}{2}\sum_{i=1}^n
    \frac{\bigl[m_{B,i} - \mathcal{M}
    - 5\log_{10}\tilde d_L(z_i;\Omega_M,\Omega_\Lambda)\bigr]^2}
    {\sigma_i^2}}_{\displaystyle\chi^2(\boldsymbol\theta)/2}.
\end{equation}
Maximising~\eqref{eq:loglik} over $\boldsymbol\theta$ minimises
$\chi^2(\boldsymbol\theta)$.

\subsection{MLE Equals WLS Under Gaussian Errors}

For fixed $\sigma_i$ (i.e., ignoring the dependence of $\sigma_{\rm int}$
on $\boldsymbol\theta$), maximising~\eqref{eq:loglik} over
$(\mathcal{M},\Omega_M,\Omega_\Lambda)$ requires minimising
\begin{equation}\label{eq:chi2}
  \chi^2(\mathcal{M},\Omega_M,\Omega_\Lambda) =
  \sum_{i=1}^n
  \frac{\bigl[m_{B,i} - \mathcal{M}
    - 5\log_{10}\tilde d_L(z_i;\Omega_M,\Omega_\Lambda)\bigr]^2}{\sigma_i^2}.
\end{equation}

Unlike the photoelectric case---where OLS had a closed-form solution---
minimising \eqref{eq:chi2} over $(\Omega_M, \Omega_\Lambda)$ requires
numerical optimisation because $\tilde d_L(z;\Omega_M,\Omega_\Lambda)$
involves a numerical integral rather than an algebraic formula.  Both
groups used grid search over a fine $(\Omega_M, \Omega_\Lambda)$ lattice,
evaluating $\chi^2$ at each grid point via numerical quadrature.
The estimator is still the minimiser:

\begin{equation}\label{eq:mle_wls}
  \boxed{(\hat{\mathcal{M}}_{\rm MLE},\,\hat\Omega_{M,\rm MLE},\,
    \hat\Omega_{\Lambda,\rm MLE})
  = \operatornamewithlimits{arg\,min}_{(\mathcal{M},\Omega_M,\Omega_\Lambda)}
    \chi^2(\mathcal{M},\Omega_M,\Omega_\Lambda).}
\end{equation}

\subsection{Profiling Out the Nuisance Parameter}

$\mathcal{M}$ can be eliminated analytically at any fixed
$(\Omega_M,\Omega_\Lambda)$, because it enters~\eqref{eq:chi2} linearly:
\begin{equation}\label{eq:Mhat}
  \hat{\mathcal{M}} =
  \frac{\displaystyle\sum_i \frac{m_{B,i} - 5\log_{10}\tilde d_L(z_i)}
    {\sigma_i^2}}
  {\displaystyle\sum_i \frac{1}{\sigma_i^2}}.
\end{equation}
Substituting~\eqref{eq:Mhat} into $\chi^2$ gives the
\emph{profile chi-squared}
\begin{equation}\label{eq:chi2_profile}
  \tilde\chi^2(\Omega_M,\Omega_\Lambda) \equiv
  \chi^2(\hat{\mathcal{M}},\Omega_M,\Omega_\Lambda),
\end{equation}
which depends only on the two shape parameters and can be minimised on a
two-dimensional grid.  In practice both groups also marginalised or
profiled over intrinsic scatter $\sigma_{\rm int}$.

\subsection{Predicted Distance Moduli for Representative Cosmological Models}

Table~\ref{tab:mu_predictions} computes the theoretical distance modulus
$\mu^{\rm th}(z;\Omega_M,\Omega_\Lambda,H_0)$ for four cosmological models
at representative redshifts, using $H_0 = 70$\,km\,s$^{-1}$\,Mpc$^{-1}$.
The model predictions diverge most strongly at $z \approx 0.3$--$0.8$.

\begin{table}[ht]
\centering
\caption{Theoretical distance modulus $\mu = 5\log_{10}(d_L/\text{Mpc})+25$
for $H_0 = 70$\,km\,s$^{-1}$\,Mpc$^{-1}$ under three cosmological
models, computed by numerical integration of equation~\eqref{eq:dl}.
$\Delta\mu_\Lambda$ is the excess distance modulus of flat $\Lambda$CDM over
the open matter-only model, and $\Delta\mu_{\rm EdS}$ its excess over
Einstein--de Sitter.  Positive values mean fainter supernovae.}
\label{tab:mu_predictions}
\medskip
\begin{tabular}{lcccrr}
\toprule
 & \multicolumn{3}{c}{$\mu^{\rm th}$ (mag)} & & \\
\cmidrule{2-4}
$z$ & Flat $\Lambda$CDM & EdS & Open &
  $\Delta\mu_\Lambda$ & $\Delta\mu_{\rm EdS}$ \\
  & $(0.3,\,0.7)$ & $(1.0,\,0.0)$ & $(0.3,\,0.0)$ &
  ($\Lambda-$open) & ($\Lambda-$EdS) \\
\midrule
0.10 & 38.32 & 38.21 & 38.25 & $+0.07$ & $+0.11$ \\
0.25 & 40.50 & 40.27 & 40.36 & $+0.14$ & $+0.24$ \\
0.50 & 42.26 & 41.86 & 42.05 & $+0.22$ & $+0.40$ \\
0.75 & 43.33 & 42.82 & 43.08 & $+0.25$ & $+0.51$ \\
1.00 & 44.10 & 43.50 & 43.84 & $+0.26$ & $+0.60$ \\
\bottomrule
\end{tabular}
\end{table}

\noindent At $z \approx 0.5$ flat $\Lambda$CDM predicts supernovae
$0.22$~mag fainter than the open matter-only model and $0.40$~mag fainter
than Einstein--de Sitter.  The total per-supernova uncertainty at that
redshift was about $0.22$~mag, of which $0.15$ is intrinsic dispersion after
standardisation and the rest photometric and $K$-correction error.  Each
event therefore carries about one standard deviation of signal against the
open model.

Naive stacking would suggest that 40 such events settle the question at
$\sqrt{40} \approx 6\sigma$.  Both teams claimed about $3\sigma$.  The gap is
the subject of Section~\ref{sec:fisher}.  Stacking assumes that the quantity
of interest is a single mean offset at fixed $\Omega_M$.  The teams estimated
a two-dimensional parameter.  Marginalising over $\Omega_M$ along the
degeneracy direction costs most of the apparent significance.
Systematic errors, which do not fall as $1/\sqrt{n}$, consumed more.

The excess faintness found in both analyses, $\Delta\mu_{\rm obs} \approx
+0.25$~mag at $z\approx0.5$ (Section~\ref{sec:results}), exceeds the
$\Lambda$CDM prediction over the open model and falls well short of the
prediction over Einstein--de Sitter.

%======================================================================
\section{Fisher Information and Confidence Ellipses}
\label{sec:fisher}

\subsection{The Fisher Information Matrix}

The precision with which $(\Omega_M, \Omega_\Lambda)$ can be estimated
is governed by the Fisher information matrix
$\mathcal{I}(\boldsymbol\theta)$.  For the model~\eqref{eq:meas2} with
known $\sigma_i$:
\begin{equation}\label{eq:fisher}
  \mathcal{I}_{jk}(\boldsymbol\theta) =
  \sum_{i=1}^n \frac{1}{\sigma_i^2}
  \frac{\partial\mu^{\rm th}(z_i;\boldsymbol\theta)}{\partial\theta_j}
  \frac{\partial\mu^{\rm th}(z_i;\boldsymbol\theta)}{\partial\theta_k},
\end{equation}
where the sum runs over all supernovae.  The Cram\'er-Rao bound states
that any unbiased estimator $\tilde{\boldsymbol\theta}$ satisfies
$\operatorname{Var}(\tilde\theta_j) \geq [\mathcal{I}^{-1}]_{jj}$.

For the shape parameters $(\Omega_M, \Omega_\Lambda)$ the relevant
derivatives are
\begin{equation}\label{eq:deriv}
  \frac{\partial\mu}{\partial\Omega_M} =
  \frac{5}{\ln 10}\,\frac{1}{\tilde d_L}\,
  \frac{\partial\tilde d_L}{\partial\Omega_M}, \qquad
  \frac{\partial\mu}{\partial\Omega_\Lambda} =
  \frac{5}{\ln 10}\,\frac{1}{\tilde d_L}\,
  \frac{\partial\tilde d_L}{\partial\Omega_\Lambda}.
\end{equation}
These are computed numerically.  $\Omega_M$ slows the expansion and reduces
$\mu$.  $\Omega_\Lambda$ accelerates the expansion and raises $\mu$.  The two
derivatives have opposite signs and comparable magnitudes over the whole
observed range, so the parameters are nearly degenerate.  At $z = 0.5$ and the
best-fit parameters,
\begin{equation}\label{eq:degeneracy}
  \frac{\partial\mu/\partial\Omega_M}{\partial\mu/\partial\Omega_\Lambda}
  \approx -1.1,
\end{equation}
falling to $-0.96$ at $z = 0.4$ and rising to $-1.24$ at $z = 0.6$.  Weighted
across an SCP-like sample the effective ratio is $-1.5$.  Raising $\Omega_M$ by
one unit and $\Omega_\Lambda$ by about $1.5$ units leaves $\mu(z)$ nearly
unchanged.

The Fisher matrix accordingly has one large and one small eigenvalue in the
$(\Omega_M,\Omega_\Lambda)$ block.  The well-measured combination is
\begin{equation}\label{eq:constraint_dir}
  0.8\,\Omega_M - 0.6\,\Omega_\Lambda,
\end{equation}
which \citet{Perlmutter1999} report as $-0.2 \pm 0.1$.  The orthogonal
combination $0.6\,\Omega_M + 0.8\,\Omega_\Lambda$ is nearly unconstrained.
The confidence ellipses are narrow along \eqref{eq:constraint_dir} and
elongated along the orthogonal direction.  \citet{GoobarPerlmutter1995}
derived this geometry before the data existed and designed the survey around
it.

\subsection{A Fisher Forecast for the 1998 Designs}
\label{sec:forecast}

Table~\ref{tab:forecast} evaluates \eqref{eq:fisher} for four designs.  Each
uses 18 low-$z$ supernovae spread over $z\in[0.02,0.10]$ with
$\sigma = 0.17$~mag, the intrinsic $0.15$ combined with the peculiar-velocity
term of Section~\ref{sec:data}, and differs only in the high-$z$ sample.  I
profile over
$\mathcal{M}$ and report standard errors on the density parameters.  The
calculation uses the redshift distributions and quoted uncertainties, not the
individual magnitudes.

\begin{table}[ht]
\centering
\caption{Fisher-information standard errors for four high-$z$ designs, each
paired with 18 low-$z$ supernovae at $\sigma = 0.17$~mag and profiled over
$\mathcal{M}$.  High-$z$ errors are $\sigma = 0.22$~mag except the HZT-like
row, which uses $0.25$~mag.  The last column is the standard error on the
combination \eqref{eq:constraint_dir}.}
\label{tab:forecast}
\medskip
\begin{tabular}{lccccc}
\toprule
High-$z$ design & $n$ & $z$ range & SE($\Omega_M$) & SE($\Omega_\Lambda$)
  & SE($0.8\Omega_M-0.6\Omega_\Lambda$) \\
\midrule
SCP-like            & 42 & 0.18--0.83 & 0.44 & 0.65 & 0.10 \\
Truncated           & 42 & 0.18--0.62 & 0.81 & 0.98 & 0.13 \\
Bunched             & 42 & 0.35--0.55 & 1.30 & 1.50 & 0.12 \\
HZT-like            & 16 & 0.16--0.62 & 1.32 & 1.56 & 0.19 \\
\bottomrule
\end{tabular}
\end{table}

The individual density parameters are not measured.  The SCP-like design gives
SE$(\Omega_M) = 0.44$ and SE$(\Omega_\Lambda) = 0.65$, so neither parameter is
determined to better than a factor of two.  The published $\pm0.09$ on
$\Omega_M$ is not a marginal standard error.  It measures precision under the
flatness constraint $\Omega_M + \Omega_\Lambda = 1$, which the same
calculation gives as $0.070$.

One combination is measured well.  The forecast standard error on
$0.8\Omega_M - 0.6\Omega_\Lambda$ is $0.10$, matching the $\pm0.1$ that
\citet{Perlmutter1999} report.  The agreement of a forecast built from
redshifts and error bars alone with the published error bar confirms that the
design fixed the 1998 constraint, and the details of the photometry did
little to it.

Redshift range buys precision on the individual parameters and almost none on
the combination.  Truncating the sample at $z = 0.62$ nearly doubles
SE$(\Omega_M)$ and widens the combination only from $0.10$ to $0.13$.  Bunching
all 42 events at $z \approx 0.45$ triples SE$(\Omega_M)$ and leaves the
combination at $0.12$.  The events beyond $z = 0.6$ rotate the ellipse.  The
rotation separates $\Omega_M$ from $\Omega_\Lambda$.

\subsection{The Value of Wide Redshift Coverage}

\citet{GoobarPerlmutter1995} showed that separating $\Omega_M$ from
$\Omega_\Lambda$ requires supernovae at two well-separated redshifts.  A sample
restricted to $z\approx 0.5$ traces a near-degenerate path in the
$(\Omega_M,\Omega_\Lambda)$ plane.  Adding $z\approx 0.8$--1.0 events rotates
the ellipse, as the second and third rows of Table~\ref{tab:forecast} show.

The SCP's highest-redshift event, SN~1997ap at $z = 0.83$, carried a large
share of that rotation.  The HZT's highest, SN~1997ck at $z = 0.97$, carried
more still, and lacked a confirming spectrum; \citet{Riess1998} reported their
conclusions with and without it.

\subsection{Numerical Illustration of the Confidence Ellipses}

\citet{Perlmutter1999} reported their confidence contours directly
as $\Delta\chi^2 = \tilde\chi^2 - \tilde\chi^2_{\rm min}$ contours.
The 68\% and 95\% confidence ellipses for one free parameter correspond
to $\Delta\chi^2 = 1$ and $4$; for two parameters, to
$\Delta\chi^2 = 2.30$ and $5.99$.  In the published confidence plots
(Figure~2 of \citealt{Perlmutter1999}; Figure~7 of \citealt{Riess1998}) the
best-fit region lies in the quadrant $\Omega_\Lambda > 0$.

The ellipse is elongated, and the line $\Omega_\Lambda = 0$ runs nearly along
the elongation, so the distance between them depends on how it is measured.
Measured by the marginal standard error on $\Omega_\Lambda$ alone, the
SCP-like forecast of Table~\ref{tab:forecast} puts $\Omega_\Lambda = 0.72$
only $1.1\sigma$ from zero.  The physical restriction $\Omega_M \geq 0$ does
the remaining work.  Setting $\Omega_\Lambda = 0$ and minimising over
$\Omega_M \geq 0$ raises $\chi^2$ by $3.8$ in the same forecast, or
$2.0\sigma$.  The published $3\sigma$ figures come from the data and the full
error budget, not from this forecast.  The forecast locates their strength.
No non-negative $\Omega_M$ reproduces the observed curvature with
$\Omega_\Lambda = 0$, and that is what excludes a decelerating universe.  The
size of $\hat\Omega_\Lambda$ on its own does not.

Under flatness the picture changes.  The same forecast gives
SE$(\Omega_M) = 0.070$, so Einstein--de Sitter ($\Omega_M = 1$) sits $10\sigma$
from the best fit.  The strength of the 1998 result rests on the joint use of
the supernova data and a prior---either $\Omega_M \geq 0$, or flatness from the
CMB.  Section~\ref{sec:limits} returns to what that dependence does and does
not license.

%======================================================================
\section{The $\mathcal{M}$ Identification Problem}
\label{sec:nuisance}

\subsection{The Contact-EMF Analogy}

In Millikan's photoelectric regression the contact electromotive force $v_c$
between the photoemitting metal and the collector entered the raw readings as
a pure intercept shift:
\[
  V_{0i}^{\rm obs} = \underbrace{-W/e - v_c}_{\alpha}
    + (h/e)\,\nu_i + \varepsilon_i.
\]
The gradient of the $V_0$--$\nu$ relation identifies the slope $h/e$.  Its
level does not.  Millikan measured $v_c$ continuously with a Kelvin
double-balance and found it constant over a run, so it shifted all
observations equally and left $\hat\beta$ unbiased.

In the supernova problem $\mathcal{M} = M_B + 5\log_{10}(c/H_0\,\text{Mpc})
+ 25$ plays the role of $v_c$.  The curvature of the $m_B$--$z$ relation
carries the cosmology.  Its level does not.  Neither $M_B$ nor $H_0$ can be
determined from the Hubble diagram alone, because a change in $H_0$ can be
offset exactly by a shift in $M_B$.  The shape parameters escape that
degeneracy.

The parallel with Millikan ends at that point, and the difference is the
harder half of the problem.  Millikan's equation was linear and its slope was
a single number, so removing $v_c$ left $h/e$ fully identified.  The Friedmann
model has two shape parameters, and removing $\mathcal{M}$ leaves them tangled
with each other.  A second degeneracy therefore survives the first: the
supernovae fix $0.8\Omega_M - 0.6\Omega_\Lambda$ and leave the orthogonal
combination open (Section~\ref{sec:fisher}).  The contact EMF was an
apparatus nuisance that better instruments could remove.  This one is a
property of the geometry, and no supernova measurement removes it.

\subsection{Constraining $\mathcal{M}$ from the Low-$z$ Sample}

Both teams' solution was identical in spirit to Millikan's: constrain
the nuisance parameter from an independent source.  At low $z$,
$d_L \approx cz/H_0$, so $m_{B,i} \approx \mathcal{M} +
5\log_{10}z_i + C$ (a constant).  The Cal\'an/Tololo low-$z$ sample
calibrates $\mathcal{M}$ through Hubble's law~\eqref{eq:hubble}, with
uncertainties from peculiar velocities averaged away over the 18 events the
SCP used.  Once $\mathcal{M}$ is anchored by the low-$z$ data, the
high-$z$ data constrain $(\Omega_M,\Omega_\Lambda)$ through the
curvature of $\mu(z)$.

\subsection{Separate Identification of $H_0$ and $M_B$}

Even with the low-$z$ anchor,  $H_0$ and $M_B$ are individually
identified only if $M_B$ is calibrated by an external distance ladder.
The primary constraint came from Cepheid variable distances measured
with the Hubble Space Telescope Key Project, which established
$H_0 = 72 \pm 8$\,km\,s$^{-1}$\,Mpc$^{-1}$ (circa 1998).  Given $H_0$,
the Cal\'an/Tololo low-$z$ sample then determines $M_B \approx -19.3$~mag.

$M_B$, equivalently $H_0$, is a nuisance parameter for the cosmological
inference.  Its uncertainty propagates into $\mathcal{M}$ but not into
$(\Omega_M,\Omega_\Lambda)$.  A 10\% error in $H_0$ shifts the Hubble diagram
by $\approx 0.22$~mag and leaves the best-fit $\Omega_M$ and $\Omega_\Lambda$
unchanged.

%======================================================================
\section{Results: The Deceleration Parameter and Cosmic Acceleration}
\label{sec:results}

\subsection{Best-Fit Parameters}

Table~\ref{tab:results} summarises the principal parameter estimates from
both groups, using their primary standardisation methods.

\begin{table}[ht]
\centering
\caption{Cosmological parameter estimates from the two Nobel-Prize groups.
Every published row imposes flatness, $\Omega_M + \Omega_\Lambda = 1$;
without it neither parameter is separately determined, as the last row
records.  Statistical uncertainties only.  The deceleration parameter is
$q_0 = \Omega_M/2 - \Omega_\Lambda$.}
\label{tab:results}
\medskip
\begin{tabular}{p{4.2cm}cccr}
\toprule
Analysis & $\hat\Omega_M$ & $\hat\Omega_\Lambda$ & $\hat q_0$ &
  Ref.\ \\
\midrule
Perlmutter et al.\ (SCP), flat   & $0.28\pm0.09$ & $0.72\pm0.09$
  & $-0.58$ & \citet{Perlmutter1999} \\
Riess et al.\ (HZT), MLCS, flat  & $0.24\pm0.10$ & $0.76\pm0.10$
  & $-0.64$ & \citet{Riess1998} \\
Riess et al.\ (HZT), template, flat & $0.20\pm0.10$ & $0.80\pm0.10$
  & $-0.70$ & \citet{Riess1998} \\
\midrule
\multicolumn{5}{l}{\emph{Standard errors without the flatness constraint
  (forecast):} SE$(\Omega_M)=0.44$, SE$(\Omega_\Lambda)=0.65$
  (Table~\ref{tab:forecast})} \\
\midrule
Planck 2018 (CMB+BAO, flat)      & $0.315\pm0.007$ & $0.685\pm0.007$
  & $-0.528$ & \citet{Planck2020} \\
\bottomrule
\end{tabular}
\end{table}

Both groups found $\hat\Omega_\Lambda$ between $0.72$ and $0.80$ from
independent datasets and different standardisation pipelines, each about one
standard error wide.  Section~\ref{sec:twogroups} compares them.  Every
published estimate in the table assumes flatness.  The supernova data alone do
not deliver these numbers.

\subsection{The Deceleration Parameter}

From the second Friedmann equation,
\begin{equation}\label{eq:q0}
  q_0 \equiv -\frac{\ddot a\,a}{\dot a^2}\bigg|_{t=t_0}
  = \frac{\Omega_M}{2} - \Omega_\Lambda.
\end{equation}
For the best-fit flat model ($\hat\Omega_M = 0.28$,
$\hat\Omega_\Lambda = 0.72$):
\begin{equation}\label{eq:q0_num}
  \hat q_0 = \frac{0.28}{2} - 0.72 = 0.14 - 0.72 = \mathbf{-0.58}.
\end{equation}
The universe is accelerating.  The transition from deceleration to
acceleration occurred at the redshift $z_t$ satisfying
$\Omega_M(1+z_t)^3 = 2\Omega_\Lambda$, which gives
\begin{equation}\label{eq:zt}
  z_t = \left(\frac{2\Omega_\Lambda}{\Omega_M}\right)^{1/3} - 1
  \approx \left(\frac{1.44}{0.28}\right)^{1/3} - 1
  \approx (5.14)^{1/3} - 1 \approx 0.73.
\end{equation}
At $z \lesssim 0.73$, dark energy dominates and expansion accelerates;
at $z \gtrsim 0.73$, matter dominates and expansion decelerates.
\citet{Riess2004} later found supernovae at $z > 1$ showing the predicted
past deceleration.  Section~\ref{sec:prediction} treats that observation and
the transition redshift they report.

\subsection{The Residual Hubble Diagram}

The residual Hubble diagram is the difference between observed distance
moduli and those predicted by the best-fit matter-only open model.  Define
\begin{equation}\label{eq:Deltamu}
  \Delta\mu_{\rm obs,i} \equiv
  \mu_{\rm obs,i} - \mu^{\rm th}(z_i;\Omega_M=0.3, \Omega_\Lambda=0).
\end{equation}
For the best-fit flat $\Lambda$CDM model,
$\Delta\mu^{\rm th}(z=0.5) \approx +0.21$~mag (Table~\ref{tab:mu_predictions}).

From both datasets, the weighted mean residual for high-$z$ supernovae
is
\begin{equation}\label{eq:Deltamu_obs}
  \overline{\Delta\mu}_{\rm obs}(z\approx0.3\text{-}0.7)
  \approx +0.25 \pm 0.06~\text{mag}.
\end{equation}
Taken as a single mean offset, the excess sits $4\sigma$ from
$\Delta\mu = 0$, and it matches the flat $\Lambda$CDM prediction to within
$1\sigma$.  That $4\sigma$ overstates the evidence for acceleration, for the
reason given in Section~\ref{sec:fisher}: the comparison holds $\Omega_M$
fixed at $0.3$, and the published significances marginalise over it.

The supernovae are fainter, hence further away, than expected.  The universe
expanded more between then and now than a decelerating universe would have.

\subsection{Statistical Significance}

For Perlmutter et al.\ (1999):
\begin{itemize}
  \item $\chi^2_{\rm min}/\nu \approx 1.03$ for the best-fit flat
    model (42 SNe, 2 free parameters after profiling $\mathcal{M}$, so
    $\nu = 39$ degrees of freedom).
  \item The no-$\Lambda$ hypothesis ($\Omega_\Lambda = 0$) gives
    $\Delta\chi^2 \approx 10.8$ (1 d.o.f.), corresponding to
    $p < 0.001$, i.e., rejected at $3.3\sigma$.
  \item Profiling over both parameters for the constraint
    $q_0 \geq 0$ (decelerating universe) gives rejection at
    $>99.9$\% confidence.
\end{itemize}
For Riess et al.\ (1998), the MLCS method gave:
\begin{itemize}
  \item $q_0 < 0$ at $3.0\sigma$ using MLCS and $2.8\sigma$ using
    template fitting.  These figures come from a secondary source and should
    be checked against \citet{Riess1998}.
  \item The two methods agreed across independent analysis pipelines, so the
    result is not an artifact of the MLCS standardisation.
\end{itemize}

%======================================================================
\section{Comparing the Two Teams}
\label{sec:twogroups}

Table~\ref{tab:comparison_teams} summarises the differences and agreements
between the two groups.

\begin{table}[ht]
\centering
\caption{Comparison of the two groups' datasets, methods, and
principal results.  Both teams reached the same conclusion from
independent data and different analysis pipelines.}
\label{tab:comparison_teams}
\medskip
\begin{tabular}{p{3.8cm}p{4.5cm}p{4.5cm}}
\toprule
Feature & High-Z Team \citep{Riess1998} &
  Supernova Cosmology Project \citep{Perlmutter1999} \\
\midrule
High-$z$ sample size & 16 SNe at $z = 0.16$--$0.62$ &
  42 SNe at $z = 0.18$--$0.83$ \\
Low-$z$ sample     & 34 SNe (literature) &
  18 SNe (Cal\'an/Tololo) \\
Standardisation    & MLCS (multi-colour light curve
  templates; $\Delta$ parameter) &
  Stretch factor ($s$; time-axis scaling of single-colour template) \\
Colour correction  & $E(B{-}V)$ from multi-colour
  light curve fit &
  $E(B{-}V)$ from $B{-}V$ at peak \\
$K$-correction    & \citet{Nugent2002} spectral templates &
  \citet{Kim1996} colour-matching \\
Intrinsic scatter  & $\sigma_{\rm int} = 0.14$~mag &
  $\sigma_{\rm int} = 0.13$~mag \\
Primary result (flat) & $\hat\Omega_M = 0.24 \pm 0.10$ &
  $\hat\Omega_M = 0.28 \pm 0.09$ \\
$q_0$ significance & $q_0<0$ at $3.0\sigma$ (MLCS) &
  $q_0<0$ at $>3\sigma$ \\
\bottomrule
\end{tabular}
\end{table}

The teams used different telescopes, observing strategies, standardisation
algorithms, and statistical analysis packages.  Their best-fit
$\hat\Omega_M$, $0.28$ and $0.24$, differ by less than half a standard error.

The agreement tests less than it appears to in one respect and more in
another.  Both numbers assume flatness, and under flatness the supernovae
measure a single parameter, so two teams agreeing on one number with
$\pm0.10$ error bars is a weak test.  Against that, the two pipelines shared
almost nothing.  MLCS and stretch parameterise the light curve differently,
the $K$-corrections came from different spectral templates, the colour
corrections used different estimators, and the high-$z$ samples overlapped
barely at all.  A systematic error in any one of those steps would have to
have been reproduced by an independent implementation to survive.  Replication
across methods tests the standardisation, not the statistics.

%======================================================================
\section{Epistemic Scepticism and Robustness Tests}
\label{sec:reluctance}

\subsection{The Initial Reaction}

Millikan's 1916 confirmation of Einstein's photoelectric equation was
greeted with scepticism about the physical interpretation, even as the
equation itself was accepted.  The supernova result of 1998--1999
provoked a structurally similar reaction: many physicists accepted the
statistical result but questioned whether it truly implied $\Omega_\Lambda
> 0$, proposing instead that systematic errors mimicked the
acceleration signal.

The principal alternative explanations advanced were:
\begin{enumerate}
  \item \textbf{Dust dimming.}  If a smooth screen of ``grey'' (wavelength-
    independent) dust occupied intergalactic space, it would absorb SN
    light at all redshifts preferentially at high $z$, making distant SNe
    appear fainter without the colour reddening that normally betrays dust.
    Both teams addressed this by checking the colour residuals of their
    SNe: if grey dust were responsible, the $B-V$ colour excess would
    scale with the magnitude offset, but no such correlation was found at
    the level required to explain the signal.

  \item \textbf{Luminosity evolution.}  If Type Ia supernovae at $z \sim
    0.5$ were intrinsically $\approx 0.25$~mag fainter than their
    low-$z$ counterparts (due to differences in progenitor age, composition,
    or environment), the signal could be entirely spurious.  The teams
    tested this by comparing the Phillips $\Delta m_{15}$ distributions
    and colour distributions of high-$z$ and low-$z$ SNe; no systematic
    difference was found.  \citet{Riess2004} later supplied the
    decisive test with supernovae at $z > 1$.  Section~\ref{sec:prediction}
    gives it.

  \item \textbf{Malmquist bias.}  At high $z$, magnitude-limited surveys
    preferentially detect intrinsically bright events, biasing $\hat m_B$
    brighter.  Both teams corrected for this with detailed simulations; the
    bias was found to be negligible compared to the acceleration signal for
    their survey depths.
\end{enumerate}
The acceleration signal survived all of these checks.

\subsection{CMB and Large-Scale Structure Evidence}

Independent corroboration came from the cosmic microwave background.
The BOOMERANG balloon experiment \citeyearpar{Boomerang2000} detected
the first acoustic peak of the CMB temperature power spectrum and
confirmed that the total energy density satisfies $\Omega_{\rm tot}
\approx 1$ (flat universe).  Combined with Hubble-constant and
large-scale-structure constraints that independently fix $\Omega_M \approx
0.3$, this implies $\Omega_\Lambda \approx 0.7$, consistent with the supernova
result within the errors of both, which were then of order $10\%$.

The corroboration matters more than the agreement.  The supernova constraint
runs along one direction of the $(\Omega_M,\Omega_\Lambda)$ plane and the CMB
constraint along a nearly orthogonal one, so the two together determine both
parameters where neither does alone.  The flatness that the published
supernova results assume is exactly what the CMB supplies.

By 2020 Planck \citep{Planck2020} gave $\Omega_\Lambda = 0.685 \pm 0.007$,
about $98$ standard errors from zero, and confirmed the 1998--1999 estimates
to sub-percent precision.

%======================================================================
\section{The Prediction at $z > 1$}
\label{sec:prediction}

The 1998 fits estimated two parameters from the data they fitted.  The test
described here estimated nothing.

Take the flat $\Lambda$CDM model with $\Omega_M = 0.28$ and
$\Omega_\Lambda = 0.72$, fitted to supernovae at $z < 0.85$, and ask what it
says about $z = 1.5$.  Referred to Einstein--de Sitter, it predicts an excess
faintness of $+0.41$~mag at $z = 0.5$ and only $+0.74$~mag at $z = 1.5$.  The
curve flattens, because dark energy dominates late and matter dominates early.
Now take the leading systematic explanation of the same 1998 data.  Grey dust
and luminosity evolution both dim distant supernovae by amounts that grow with
distance.  Calibrated to reproduce the observed $+0.41$~mag at $z = 0.5$, a
dimming proportional to $z$ predicts $+1.24$~mag at $z = 1.5$ and keeps
climbing.

\begin{table}[ht]
\centering
\caption{Excess faintness relative to Einstein--de Sitter under the flat
$\Lambda$CDM model fitted in 1998, and under a dimming proportional to $z$
calibrated to match it at $z = 0.5$.  The last column is the difference the
two models predict, in magnitudes.  Per-supernova errors at these redshifts
were $0.2$--$0.3$~mag.}
\label{tab:prediction}
\medskip
\begin{tabular}{lrrr}
\toprule
$z$ & $\Lambda$CDM & Dimming $\propto z$ & Difference \\
\midrule
0.5 & $+0.41$ & $+0.41$ & $0.00$ \\
1.0 & $+0.62$ & $+0.83$ & $-0.20$ \\
1.5 & $+0.74$ & $+1.24$ & $-0.50$ \\
\bottomrule
\end{tabular}
\end{table}

The two models separate by half a magnitude at $z = 1.5$, two to three
per-supernova standard deviations, which a handful of events settles.
Referred instead to an empty universe the prediction is sharper still.
$\Lambda$CDM puts supernovae beyond $z \approx 1.4$ \emph{brighter} than that
benchmark, the signature of the deceleration that preceded the dark-energy
era.  No model in which distant supernovae are dimmed by dust or by evolution
produces a change of sign.

\citet{Riess2004} collected the supernovae at $z > 1$ with the Hubble Space
Telescope and found the turnover.  The 1998 parameters had no freedom left to
accommodate the result.  That observation is the out-of-sample test of the
model.  The $3\sigma$ of 1998 measured parameters within it.

One clarification about the transition redshift.  The flat model fitted in
1998 implies $z_t = (2\Omega_\Lambda/\Omega_M)^{1/3} - 1 = 0.73$, and the
Planck parameters imply $0.63$.  The figure $z_t \approx 0.46$ that
\citet{Riess2004} report comes from a kinematic fit in which $q(z)$ is linear
in $z$, not from $\Lambda$CDM, and the two numbers are not comparable.  The
qualitative prediction---that deceleration preceded acceleration, and that the
turnover falls near $z \sim 0.5$--$0.8$---is what the observation confirmed.

%======================================================================
\section{What the Supernovae Could Not Settle}
\label{sec:limits}

The robustness tests of Section~\ref{sec:reluctance} answered questions about
the data.  A separate question concerns the model, and the supernovae do not
answer it.

Type~Ia supernovae measure one function, the luminosity distance $d_L(z)$.
Every quantity in this paper is a functional of that one curve.  Two theories
that deliver the same $d_L(z)$ over $z \in [0, 1]$ are observationally
identical to these data, whatever else they assert.  Three classes of theory
do so.  A cosmological constant with $w = -1$ is one.  A rolling scalar field
with $w(z)$ within a few percent of $-1$ over the observed range is a second,
and Appendix~\ref{app:eos} records that current data bound $w$ only to about
$0.03$.  Modified theories of gravity that reproduce the same expansion
history without any dark-energy fluid are a third.  The 1998 data select the
functional form of $d_L(z)$.  They do not select the physics that generates
it.

The inference to acceleration itself rests on a prior.
Section~\ref{sec:fisher} showed that the supernovae measure
$0.8\Omega_M - 0.6\Omega_\Lambda$ to $\pm0.1$ and the orthogonal combination
hardly at all.  Ruling out $\Omega_\Lambda = 0$ requires either
$\Omega_M \geq 0$, which is a physical restriction rather than a measurement,
or flatness, which the CMB supplies.  Neither prior is controversial.  Both
are external to the supernova data, and the published error bars do not
survive their removal.

Distinguishing $\Lambda$ from a dynamical field requires a different
observable.  \citet{Riess2004} and the later surveys pursued one of them, the
effect of $w$ on $d_L(z)$ at higher redshift.  The growth rate of large-scale
structure supplies the other.  Modified gravity and dark energy predict
different growth rates even when they produce identical expansion histories,
so growth separates them where distances cannot.  Theory names the
discriminating observable.  No one would measure growth rates to test dark
energy without a theory predicting that the two differ.

The mechanism remains unexplained.  Quantum field theory predicts a vacuum
energy density some $10^{120}$ times the observed $\Lambda$
\citep{Weinberg1989}.  Riess, Schmidt, and Perlmutter measured the curvature
of the Hubble diagram to a few percent.  The discrepancy between that
measurement and the theory of what it measures is the largest in physics.

%======================================================================
\section{The Physicists as Machine Learners}
\label{sec:ml}

\subsection{Overview}

The vocabulary of modern machine learning---supervised learning, hypothesis
class, loss function, inductive bias, feature engineering, transfer
learning---did not circulate in the astrophysics community of the 1990s.
The two teams' reasoning maps onto it as Einstein's 1905 reasoning does.
\citet{sargent2025sources} argues that the mapping generally holds, and that
the techniques now gathered under artificial intelligence have long histories
inside the sciences that used them.  This section states the correspondence
for the supernova case.

\subsection{The Hypothesis Class}

The Friedmann equations and the definition of luminosity distance fix the
hypothesis class for the supernova problem:
\begin{equation}\label{eq:hyp_class}
  \mathcal{H} = \bigl\{
    m_B \mapsto \mathcal{M} + 5\log_{10}\tilde d_L(z;\Omega_M,\Omega_\Lambda)
    \bigm|
    \mathcal{M}\in\mathbb{R},\;
    \Omega_M \geq 0,\;
    \Omega_\Lambda \in\mathbb{R}
  \bigr\}.
\end{equation}
This is a three-parameter family of smooth curves on the $(z,\,m_B)$ plane.
Without the Friedmann equations the space of possible $m_B(z)$ functions is
infinite-dimensional; with them it collapses to the three real numbers
$({\mathcal M}, \Omega_M, \Omega_\Lambda)$.

The theoretical restriction to $\mathcal{H}$ was not arbitrary.  It follows
deductively from three independent bodies of physics:
\begin{enumerate}
  \item General relativity (the geodesic equation in an expanding spacetime);
  \item The assumption of spatial homogeneity and isotropy (the cosmological
    principle);
  \item The assumption that pressure-free matter obeys $\rho_M \propto a^{-3}$
    and the cosmological constant is truly constant.
\end{enumerate}
Each assumption had been tested independently long before the supernova search.
The Friedmann hypothesis class was well-motivated, not curve-fitted.

\subsection{Feature Engineering: The Phillips Width--Luminosity Relation}

Raw photometry of Type~Ia supernovae involves hundreds of measurements at
different epochs and wavelengths.  \citet{Phillips1993} found that the
post-peak decline rate $\Delta m_{15}$ predicts the peak luminosity, which
reduces that high-dimensional representation to a one-dimensional
correction.  The correction cut the dispersion from $0.4$ to $0.15$~mag, a
factor of $2.7$ in the noise and a factor of $7$ in the number of supernovae
needed for a given precision.

Against the open matter-only model the signal at $z\approx0.5$ is $0.22$~mag
(Table~\ref{tab:mu_predictions}).  Detecting a $0.22$~mag offset at $3\sigma$
takes about $30$ supernovae at $\sigma = 0.4$~mag and about $4$ at
$\sigma = 0.15$~mag, before any allowance for the parameter degeneracy of
Section~\ref{sec:fisher} or for systematics.  Both teams had roughly the
number of events that the pre-Phillips dispersion would have required merely
to see the effect, and neither could have separated it from $\Omega_M$.

A data-driven learner given raw light curves would have to extract the
Phillips representation from the data alone.  Phillips did exactly that.

\subsection{The Inductive Half of the Discovery}
\label{sec:phillips_induction}

The account so far credits theory with the work.  General relativity fixed the
functional form, three parameters remained, and fifty supernovae sufficed
where a theory-free learner would have needed thousands.  Half of the
discovery does not fit that account.

$\Delta m_{15}$ came from data.  No theory of Type~Ia explosions predicted
that the post-peak decline rate tracks peak luminosity, and
\citet{Phillips1993} did not derive the relation.  He plotted decline rate
against absolute magnitude for nine nearby supernovae and fitted a line.  The
relation has since been partly rationalised---more $^{56}$Ni yields both a
brighter peak and a more opaque, slower-cooling ejecta---and it has not been
deduced.  Thirty years on, the width--luminosity relation remains an empirical
regularity whose calibration is refitted with every new survey.

The progenitor physics behind it is unsettled in a way that matters here.
Section~\ref{sec:standard_candle} records that the Chandrasekhar-mass account
no longer commands consensus.  A theorist in 1993 with the correct progenitor
model in hand could not have written down equation~\eqref{eq:phillips}, and a
theorist today still cannot.

The supernova discovery was therefore a hybrid.  A deductively specified
hypothesis class was bolted to an inductively discovered feature, and neither
half would have sufficed.  Without Friedmann there is no $d_L(z)$ to fit and
no meaning for the fitted numbers.  Without Phillips the scatter is twice the
signal and no amount of theory recovers it.  The one step in the episode where
a data-driven procedure did the indispensable work was the one step where no
theory was available to do it instead.

The pattern has continued.  Modern supernova cosmology runs on SALT2 and
SALT3, light-curve models trained on data rather than derived from explosion
physics, and their training sets grow with each survey.  Feature extraction in
this field has become more machine-learned over time.  The hypothesis class
has stayed Friedmann's.  The division of labour that Section~\ref{sec:dnn}
describes is not a claim that induction fails; it is a claim about which part
of the problem each method can reach.

\subsection{Training Data and Generalisation}

The \emph{training data} for both analyses consisted of the low-$z$
Cal\'an/Tololo supernovae \citep{Hamuy1996}: the $\approx 29$ events
at $z = 0.01$--$0.10$ that calibrated the nuisance parameter
$\mathcal{M}$ (and, via $H_0$, the absolute magnitude $M_B$).

The high-$z$ events at $z = 0.3$--$0.8$ carried most of the cosmological
information.  They were not held out.  Both $(\Omega_M,\Omega_\Lambda)$ and
$\mathcal{M}$ were fitted to the low-$z$ and high-$z$ samples together, so
neither set is test data for the cosmology.  Section~\ref{sec:prediction}
describes the observation that was.

The standardisation did transfer out of sample.  Both teams applied the MLCS
or stretch correction learned from low-$z$ supernovae, unchanged, to the
high-$z$ events.
\begin{itemize}
  \item The HZT's MLCS templates and the coefficient $\eta = 0.91$~mag
    were learned from low-$z$ training data and then used to correct
    high-$z$ events (out-of-sample prediction).
  \item The SCP's stretch factor $s$ and correction coefficient
    $\alpha_{\rm SCP}$ were similarly trained on low-$z$ data.
\end{itemize}
The implicit assumption---that Type~Ia supernovae at $z\approx 0.5$
follow the same Phillips relation as nearby events---is a
\emph{transfer learning} assumption.  Its validity was tested through
the spectroscopic comparison of high-$z$ and low-$z$ SN~Ia spectra:
both teams found that the optical spectra of high-$z$ events were,
within measurement uncertainty, identical to those of low-$z$ events
of the same light-curve shape.

The cross-$z$ universality of the shape--luminosity relation reduces the
effective degrees of freedom and makes the inference tractable with a small
sample.

\subsection{Model Selection: Accelerating versus Decelerating Universes}

At the time of both papers, four main cosmological models competed:
\begin{itemize}
  \item \textbf{Einstein-de Sitter ($\Omega_M = 1,\,\Omega_\Lambda = 0$):}
    the theoretically favoured flat matter-only model.
  \item \textbf{Open matter-only ($\Omega_M \approx 0.3,\,\Omega_\Lambda =
    0$):} consistent with galaxy cluster mass measurements.
  \item \textbf{Flat $\Lambda$CDM ($\Omega_M \approx 0.3,\,\Omega_\Lambda
    \approx 0.7$):} the dark-energy model.
  \item \textbf{Flat ``QCDM'' with dynamical dark energy ($w\neq -1$):}
    generalisations where $\Omega_\Lambda$ is replaced by a field with
    equation of state $w < -1/3$.
\end{itemize}
The supernova data, combined with the flat-universe CMB evidence, selected
flat $\Lambda$CDM: the best-fit $\chi^2$ under the winning model was lower by
$\Delta\chi^2 > 10$ than under Einstein--de Sitter.  Any modern criterion---AIC,
BIC, Bayesian evidence---returns the same verdict.

The verdict is narrower than it sounds.  The first three models on the list
differ only in $(\Omega_M,\Omega_\Lambda)$, so the selection among them is
parameter estimation within one hypothesis class.  The fourth differs in $w$,
and the supernova data of 1998 could not distinguish $w = -1$ from
$w = -0.8$.  The data rejected a decelerating expansion history.  They left
standing every theory that reproduces the observed $d_L(z)$, which
Section~\ref{sec:limits} enumerates.

\subsection{The Deductive-Inductive Loop}

Table~\ref{tab:loop} translates the supernova discovery into the language
of modern machine learning.

\begin{table}[ht]
\centering
\caption{The hypothetico-deductive loop in the supernova episode,
translated into the language of modern machine learning.}
\label{tab:loop}
\begin{tabular}{p{3.6cm}p{4.0cm}p{4.0cm}}
\toprule
Step & Astrophysical name & ML name \\
\midrule
Propose Friedmann $+\Lambda$ & Theoretical model & Model architecture \\
Derive $d_L(z;\Omega_M,\Omega_\Lambda)$ & Prediction from GR &
  Hypothesis class \\
Phillips relation: use $\Delta m_{15}$, not raw $m_B$ &
  Standardisation & Feature engineering \\
Cal\'an/Tololo survey & Low-$z$ calibration &
  Training data \\
High-$z$ SN Ia searches & Discovery campaigns &
  Training-data extension \\
Minimise $\chi^2(\Omega_M,\Omega_\Lambda)$ & Maximum likelihood &
  WLS training \\
$\Omega_\Lambda > 0$ at $3\sigma$ & Acceleration inferred &
  Parameter estimate \\
Deceleration seen at $z>1$ & Prediction confirmed &
  Out-of-sample test \\
Same result from SCP and HZT & Cross-group replication &
  Independent reimplementation \\
\bottomrule
\end{tabular}
\end{table}

%======================================================================
\section{A Deep Neural Network on the Supernova Data}
\label{sec:dnn}

\subsection{Two Hypothesis Classes Compared}

The Friedmann-equations hypothesis class~\eqref{eq:hyp_class} has three
free parameters after fixing $\sigma_{\rm int}$.  A deep neural
network (DNN) for the regression $z \to m_B$ is a universal function
approximator \citep{Hornik1989}: a three-layer network with 16 ReLU
neurons per layer has approximately
\[
  p = (1\cdot16+16) + (16\cdot16+16) + (16\cdot16+16) + (16\cdot1+1) = 593~\text{parameters}.
\]
For ReLU networks the pseudo-dimension grows like $p\log_2 p$, here about
$5{,}500$ \citep{Vapnik1998}.  Converting that into a required sample size
admits no single answer.  The uniform-convergence bound $O(d/\epsilon^2)$
gives $n\sim10^7$ for $\epsilon = 1\%$ and is vacuous at that size.  The
practitioner's rule $n \gtrsim 10p$ gives $n \gtrsim 6{,}000$, and I use it
below.  The SCP had 42 supernovae; the HZT had 50.

Capacity bounds of this kind are loose.  Overparameterised networks
generalise in regimes where uniform-convergence theory says they cannot
\citep{Belkin2019}, and nothing below rests on a claim that a network provably
fails.  The binding difficulty at $n = 42$ is identifiability rather than
generalisation error.

\subsection{Epoch I: Before Phillips (1993)}

Before the width--luminosity relation was known, the Type~Ia scatter
was $\approx 0.4$~mag.  A DNN applied to raw $(z, m_B)$ data would face
an $n=O(20)$ dataset with $\sigma_{\rm noise} \approx 0.4$~mag.

\paragraph{Underdetermination.}
With $n = 20$ and $p = 593$, the DNN is massively overparameterised.
Gradient descent finds a minimum-norm interpolant that passes through
all 20 training points but is unconstrained between them
\citep{Belkin2019}.

\paragraph{Loss of the discrimination.}
The cosmological signal is $\approx 0.25$~mag at $z \approx 0.5$---
smaller than the raw $0.4$~mag scatter.  Without the Phillips feature
transformation, the DNN cannot separate signal from noise at all.

\subsection{Epoch II: After Phillips (1993), Before Large Samples}

The Cal\'an/Tololo survey (1996) provides clean standardised photometry for
29 supernovae, of which the SCP used 18.  Either number is far short of what a
593-parameter network requires.

\paragraph{No residual degrees of freedom.}
The Friedmann model leaves $n - 3 = 15$ residual degrees of freedom for
$n = 18$ observations and three free parameters---enough to estimate
$\hat\sigma_{\rm int}$ and compute meaningful confidence intervals.
The DNN has no such decomposition: with 593 parameters and 18 data points,
every weight and bias is undertrained; there are no residual degrees
of freedom to estimate noise variance or standard errors.

\paragraph{Cannot identify the degeneracy direction.}
The degeneracy direction $0.8\,\Omega_M - 0.6\,\Omega_\Lambda \approx
\text{const}$ (Section~\ref{sec:fisher}) must be identified from the
\emph{curvature} of $\mu(z)$ over a wide redshift range.  A DNN trained
on 18 low-$z$ events, all bunched in the near-linear regime $z < 0.1$,
has no way to estimate this curvature: the functional form is undetermined.
The Friedmann model, by contrast, predicts the curvature analytically
from the density parameters, requiring only a small number of well-placed
high-$z$ observations to pin it down.

\subsection{Epoch III: The 1998--1999 Discovery Datasets}

The SCP's 42-SN and HZT's 50-SN samples are the operative training sets
for the discovery.

\paragraph{Overfitting in the interpolation regime.}
With $n = 42$ and $p = 593$, the DNN still operates in the interpolation
regime: it can memorise all 42 training points exactly,
but its predictions between and beyond the training redshifts are
unconstrained.  In particular, the DNN has no mechanism to recover the
physically meaningful separation $\hat q_0 = -0.58$ from the overfit
weights.

\paragraph{No Cramér-Rao bound.}
The cosmological analysis produced confidence ellipses from
$\Delta\chi^2$ contours; the Friedmann model reduces the fit to
$n - 3 = 39$ degrees of freedom.  A DNN has no
analogue: because its residuals are zero by construction (interpolation),
there is no estimate of $\sigma_{\rm int}$, no prediction interval,
and no confidence surface over $(\Omega_M, \Omega_\Lambda)$.

\paragraph{No physical interpretation of weights.}
Even if the DNN happened to learn a curve close to $\mu(z;0.28,0.72)$,
the weights of the network would not decompose into $(\hat\Omega_M,
\hat\Omega_\Lambda)$.  The connection between the learned function and
the deceleration parameter $q_0 = \Omega_M/2 - \Omega_\Lambda$ would be
opaque.

\subsection{The Modern Era: Large Supernova Compilations}

By the 2010s, supernova compilations reached several hundred events
(e.g., \citealt{Conley2011}, Joint Light-Curve Analysis).  In this
regime a DNN with appropriate architecture and regularisation could,
in principle, learn the shape of $\mu(z)$ from the data alone and
converge toward the Friedmann functional form.  It would still not identify
the function, extract the parameters, relate them to $q_0$, or extrapolate
past the training range:

\begin{itemize}
  \item Identify that the learned function is $5\log_{10}\tilde d_L(z)$
    where $\tilde d_L$ satisfies a specific integral from general
    relativity.
  \item Extract $(\Omega_M, \Omega_\Lambda)$ from the network weights
    without supplying the Friedmann equations as an external constraint.
  \item Relate the curvature of $\mu(z)$ to the deceleration parameter
    $q_0 = \Omega_M/2 - \Omega_\Lambda$ and hence to the question
    of whether the universe is accelerating.
  \item Predict the $z > 1$ deceleration observed by \citet{Riess2004}:
    the DNN trained on $z < 1$ data would extrapolate the
    accelerating trend, not the decelerating one, because it has no
    physical constraint requiring the cosmological transition.
\end{itemize}

Table~\ref{tab:sample_complexity} quantifies the contrast between the
Friedmann model and a DNN at each epoch of the supernova program.

\begin{table}[ht]
\centering
\caption{Sample-complexity comparison between the three-parameter
Friedmann model (pseudo-dimension $= 3$) and a three-layer DNN
($p = 593$ parameters, pseudo-dimension $\approx 5{,}500$).  ``Sufficient
$n$'' uses the practitioner's rule $n \gtrsim 10p \approx 6{,}000$ for the
DNN, and four for the Friedmann model, one more than its free parameters.
Sufficiency for the Friedmann model means enough residual degrees of freedom
to fit and to assess the fit; it does not mean the parameters are identified,
which requires high-$z$ coverage as well.}
\label{tab:sample_complexity}
\begin{tabular}{p{2.6cm}rp{2.4cm}p{2.4cm}p{3.8cm}}
\toprule
Epoch & $n$ available & Friedmann model & DNN
  (VC$\approx 5700$) & DNN outcome \\
\midrule
Pre-Phillips (before 1993) & $\sim 20$ raw & Insufficient$^\dagger$
  & Insufficient & Cannot detect signal (noise $> $ signal) \\
Cal\'an/ Tololo 1996 & 29 (low-$z$) & Insufficient$^\dagger$
  & Insufficient & Cannot constrain curvature of $\mu(z)$ \\
Discovery 1998 & 42--50 & Sufficient
  & Insufficient ($\times$140) & Overfits; no inference \\
JLA compilation 2011 & $\sim 500$ & Sufficient
  & Insufficient ($\times$12) & Partially learns shape; cannot extract $q_0$ \\
Modern era 2020s & $\sim 2000$ & Sufficient
  & Insufficient ($\times$3) & Recovers $\mu(z)$ shape; cannot identify $q_0$ \\
\bottomrule
\multicolumn{5}{p{14cm}}{$^\dagger$ Without high-$z$ data, $(\Omega_M,
  \Omega_\Lambda)$ are not identified: only $\mathcal{M}$ can be
  constrained from low-$z$ observations.}
\end{tabular}
\end{table}

The photoelectric episode has the same structure.  Einstein's quantum
hypothesis reduced an infinite-dimensional function-learning problem to a
two-parameter line-fitting problem.  Friedmann's equations and the
cosmological principle reduced the supernova distance-magnitude relation to a
three-parameter curve, and Phillips reduced the scatter enough to make that
curve measurable from $\approx 40$ events.  In both cases the theorists
specified the smallest useful hypothesis class; the data then constrained the
free parameters within it.

%======================================================================
\section{Conclusions}
\label{sec:conclusions}

The supernova discovery of cosmic acceleration is, at its mathematical
core, a three-parameter weighted nonlinear regression.  The free
parameters are $\mathcal{M}$ (a nuisance combination of absolute
magnitude and Hubble constant), $\Omega_M$ (matter density), and
$\Omega_\Lambda$ (cosmological constant / dark energy density).  The
measurement equation
\[
  m_{B,i} = \mathcal{M} + 5\log_{10}\tilde d_L(z_i;\Omega_M,\Omega_\Lambda)
  + \varepsilon_i
\]
is derived directly from general relativity; the Friedmann equations
are the physical hypothesis that specifies the functional form.

The regression is straightforward.  Two other tasks were not.
Standardising the supernovae by the Phillips width--luminosity relation cut
the scatter from $0.4$ to $0.15$~mag, a factor of $2.7$ in the noise and a
factor of seven in the number of events needed.  Assembling supernovae at
$z = 0.3$--$0.8$ put observations where the accelerating and decelerating
models differ by $0.2$--$0.3$~mag.  Neither task is statistical.

The estimates of both groups, under flatness, give
\begin{align*}
  \hat q_0 &= \frac{\hat\Omega_M}{2} - \hat\Omega_\Lambda
           = 0.14 - 0.72 = -0.58, \\
  z_t      &\approx 0.73, \\
  \hat\Omega_\Lambda &= 0.72 \pm 0.09~\text{(stat)}.
\end{align*}

Roughly fifty supernovae sufficed because the Friedmann equations fixed the
functional form and left three parameters.  A learner obliged to discover the
form from data would need thousands of events and would still not know what it
had found.  Theory compresses the sample size that data must supply.

The 1998 fits estimated parameters.  The model's prediction that supernovae
beyond $z \approx 1.4$ would appear brighter than an empty universe estimated
nothing, and \citet{Riess2004} confirmed it
(Section~\ref{sec:prediction}).  No systematic-error model produced that
change of sign.

The supernovae measured one function, $d_L(z)$, and one combination of two
parameters within it.  Separating $\Omega_M$ from $\Omega_\Lambda$ required the
CMB, whose constraint runs nearly orthogonal to theirs.  Distinguishing a
cosmological constant from a rolling field will require the growth rate of
structure, which no one would measure without a theory predicting that the two
differ.  Theory sets the hypothesis class.  Theory then names the measurement
that discriminates within it.

Riess, Schmidt, and Perlmutter measured the curvature of the Hubble diagram
and found $\Omega_\Lambda \approx 0.72$.  They did not explain what dark
energy is.  Quantum field theory still predicts a vacuum energy some
$10^{120}$ times larger \citep{Weinberg1989}.

%======================================================================
\appendix

\section{The Deceleration Parameter and the Second Friedmann Equation}
\label{app:q0}

The deceleration parameter $q_0$ follows from the second Friedmann equation.
With matter (pressure $p_M = 0$) and cosmological constant
($p_\Lambda = -\rho_\Lambda c^2$):
\begin{equation}\label{eq:friedmann2}
  \frac{\ddot a}{a} = -\frac{4\pi G}{3}\left(\rho + \frac{3p}{c^2}\right)
  + \frac{\Lambda}{3}
  = -\frac{4\pi G \rho_M}{3} + \frac{\Lambda}{3}.
\end{equation}
Evaluating at $t = t_0$ (today, $\rho_{M,0} = 3H_0^2\Omega_M/(8\pi G)$,
$\Lambda = 3H_0^2\Omega_\Lambda$):
\begin{equation}
  \frac{\ddot a_0}{a_0} = -\frac{H_0^2\Omega_M}{2} + H_0^2\Omega_\Lambda.
\end{equation}
Therefore
\begin{equation}\label{eq:q0_derive}
  q_0 = -\frac{\ddot a_0 a_0}{\dot a_0^2}
  = -\frac{\ddot a_0/a_0}{H_0^2}
  = \frac{\Omega_M}{2} - \Omega_\Lambda.
\end{equation}
For $\Omega_M < 2\Omega_\Lambda$ the universe is accelerating today.
The best-fit flat model ($\Omega_M = 0.28$, $\Omega_\Lambda = 0.72$)
gives $q_0 = 0.14 - 0.72 = -0.58$: the universe currently accelerates
at a rate $|q_0| = 0.58$ times the Hubble deceleration.

\section{Low-Redshift Expansion and the Hubble Diagram}
\label{app:lowz}

At low redshifts, the luminosity distance has a Taylor expansion in~$z$:
\begin{equation}\label{eq:dl_expand}
  d_L(z) = \frac{cz}{H_0}\left[1 + \tfrac{1}{2}(1-q_0)z + O(z^2)\right].
\end{equation}
The distance modulus then expands as
\begin{equation}\label{eq:mu_expand}
  \mu(z) = 5\log_{10}\!\left(\frac{cz}{H_0\,\text{Mpc}}\right) + 25
  + \frac{5(1-q_0)}{2\ln 10}\,z + O(z^2).
\end{equation}
The coefficient $5/(2\ln 10) = 1.086$, so the cosmological term is
$1.086\,(1-q_0)\,z$.  At $z = 0.1$ and $q_0 = -0.58$ it equals $0.17$~mag,
which exceeds $\sigma_{\rm int} \approx 0.15$~mag.  The term is not
negligible, and the low-$z$ sample is not free of cosmological information.

Leverage defeats the low-$z$ sample.  The
coefficient multiplies $z$, and the Cal\'an/Tololo events span only
$z \in [0.01, 0.10]$.  For 29 events spread over that range,
$\sqrt{S_{zz}} = 0.145$, so the standard error on the coefficient is
$\sigma/\sqrt{S_{zz}} = 0.17/0.145 = 1.17$, and
\begin{equation}\label{eq:seq0_lowz}
  \mathrm{SE}(q_0) = \frac{1.17}{1.086} \approx 1.1 .
\end{equation}
A standard error of $1.1$ on a parameter whose competing values are $-0.58$
and $+0.5$ settles nothing.  The low-$z$ sample constrains $\mathcal{M}$
through Hubble's law and leaves $q_0$ open.

The same expansion prices the alternative.  Moving to $z = 0.5$ multiplies the
signal by five and multiplies $\sqrt{S_{zz}}$ by more, because the leverage
grows with the spread of $z$ and not with its level.  Millikan's
frequency-range argument has the same form, and the same consideration drove
both teams' observing strategy.

\section{The Equation of State of Dark Energy}
\label{app:eos}

The cosmological constant $\Lambda$ corresponds to a fluid with equation
of state $w_\Lambda = p_\Lambda/(\rho_\Lambda c^2) = -1$: its pressure
is equal and opposite to its energy density.  A dark-energy fluid with
$w \neq -1$ (e.g., a rolling scalar field $\phi$ with potential $V(\phi)$)
also has $\ddot a > 0$ provided $w < -1/3$.

The most general measurement equation, allowing $w$ as a free parameter, is
\begin{equation}\label{eq:mu_w}
  d_L(z;\Omega_M,\Omega_{\rm de},w) = \frac{c(1+z)}{H_0}
  \int_0^z \frac{dz'}{\sqrt{\Omega_M(1+z')^3
    + \Omega_{\rm de}(1+z')^{3(1+w)}}}.
\end{equation}
The 1998--1999 analyses assumed $w = -1$.  Later surveys---SNLS, Union, JLA,
DES---relaxed it \citep{Conley2011,Sullivan2011}.  Combined supernova, CMB and
baryon-acoustic-oscillation data now give $w = -1.03 \pm 0.03$, consistent
with a cosmological constant.  That figure comes from a secondary source and
should be checked against the primary analyses.

The precision on $w$ understates how little it settles.  A rolling scalar
field with $w$ within a few percent of $-1$ over the observed redshift range
is indistinguishable from $\Lambda$ by these data, and so are modified-gravity
models tuned to reproduce the same expansion history.  Supernovae measure
$d_L(z)$.  Any theory delivering the same $d_L(z)$ is observationally
identical to any other, and the remaining uncertainty in $w$ now comes from
systematics rather than photon statistics.

\bigskip
\bibliographystyle{plainnat}
\bibliography{riess_schmidt_perlmutter_ML}

\end{document}
